% ====================================================================
% DESCRICAO-ALTO-NIVEL.TEX — Seção 4: Descrição de Alto Nível
% ====================================================================

\section{Descrição de Alto Nível}

\subsection{Visão Geral do Sistema}

O sistema de monitoramento ambiental é composto por um \textbf{microcontrolador ESP32-WROOM-32U}
que atua como núcleo central, integrando a leitura de dados ambientais, a comunicação
em rede e o controle dos atuadores (LEDs).

O fluxo geral de operação é o seguinte:

\begin{enumerate}
    \item O \textbf{sensor DHT11} coleta periodicamente dados de temperatura e umidade
    do ambiente e os envia ao ESP32 via protocolo digital single-wire.

    \item O \textbf{ESP32} processa os dados recebidos e os disponibiliza em uma
    \textbf{interface web} acessível em tempo real via Wi-Fi, sem necessidade de
    recarregar a página (atualização dinâmica).

    \item O usuário acessa a interface web e pode \textbf{alternar entre os modos de
    operação} (Desligado, Ligado e Automático) por meio de botões na página.

    \item Ao mudar o modo, o ESP32 aciona os \textbf{LEDs indicadores} (LED1, LED2 e LED3)
    para refletir visualmente o estado ativo do sistema.

    \item O \textbf{módulo MB102} fornece a alimentação regulada (3,3 V / 5 V) para
    todos os componentes do circuito.
\end{enumerate}


% --------------------------------------------------------------------
\subsection{Modos de Operação}
% --------------------------------------------------------------------

O sistema opera em três modos distintos, selecionáveis pelo usuário via interface web:

\begin{itemize}
    \item \textbf{Modo Desligado:} o sensor DHT11 e o sistema de coleta são desativados.
    O LED3 acende para indicar o estado inativo. Nenhum dado é coletado ou exibido.

    \item \textbf{Modo Ligado:} o sistema opera normalmente, aguardando ação do usuário
    para realizar e exibir leituras. O LED1 acende para indicar este modo.

    \item \textbf{Modo Automático:} o ESP32 realiza leituras contínuas e automáticas
    do DHT11, atualizando a interface web em tempo real sem intervenção do usuário.
    O LED2 acende para indicar este modo.
\end{itemize}

% --------------------------------------------------------------------
% NOTA: Dado o escopo reduzido e bem definido do projeto, não foi
% identificada a necessidade de divisão em subsistemas formais.
% O sistema é suficientemente simples para ser descrito como um
% bloco único com três periféricos (sensor, LEDs e interface web).
% Caso o projeto evolua (app mobile, log local, alertas por e-mail),
% considere decompor em subsistemas de hardware, firmware e frontend.
% --------------------------------------------------------------------
